\section{Conclusion}
\label{sec:conclusion}

This companion paper develops the recursive structural decomposition
design choice that the framework paper at \texttt{paper/} introduces
through schema-level parent-pointer fields and then defers. We argued
five points: that structural recursion has two operationally distinct
modes (productive decomposition that delivers leaf-level verifier signal,
and unproductive meta-criticism that does not); that a bounded depth
limit $\depthlim$ acts as a structural prior on the MDP transition
function and selects the productive mode without foreclosing the
unproductive one; that the framework paper's seven-verb action set
$\Aset$ extends to arbitrary recursive depth via identifier-based
targeting, with no new categorical actions and no inflation of the PUCT
branching factor; that flat serialization with depth markers lets a
frozen LLM resolve recursive tree structure through native causal
self-attention, with no graph neural network and no tree-LSTM; and that
the central depth-limit claim is falsifiable via a pre-registered
ablation grid that varies $\depthlim$ across $\{0, 3, \infty\}$ on the
framework paper's benchmarks.

The contribution is methodological. The recursive paper specifies when
and how the framework paper's schema-level recursion hooks should be
exercised, and the pre-registered ablation makes the central claim a
falsifiable empirical question. The expected ablation outcome is that
$\depthlim = 3$ outperforms both endpoints: $\depthlim = 0$ loses the
expressivity that productive decomposition needs, and $\depthlim =
\infty$ wastes sample budget on unproductive trajectories whose Bellman
targets contain no informative verifier signal. Whether this prediction
holds empirically is the question the experiment will answer, and the
framework paper's implementation is where the answer will be reported.
